%! TeX program = lualatex
\documentclass[10pt,a4paper]{article}

\usepackage[english]{babel}
\usepackage{fontspec}
\setmainfont{Times New Roman}
\newfontfamily\vest{TRY Vesterbro Regular}
\newfontfamily\vestb{TRY Vesterbro Poster}
\newfontfamily\overp{Overpass}
\usepackage{graphicx}
\graphicspath{{../../img/}}
\usepackage[svgnames]{xcolor}
\definecolor{frbblue}{RGB}{47,88,109}
\definecolor{frbgreen}{RGB}{28,85,66}
\usepackage[pdfborder={0 0 0}]{hyperref}
\usepackage{enumitem}
\usepackage[modulo]{lineno}
\usepackage{titlesec}
\titleformat{\subsection}{\color{frbgreen}\overp\large}{}{}{}
\usepackage{titling}
\setlength{\droptitle}{-12ex}
\pretitle{\begin{center}\vestb\Huge\color{frbgreen}}
\posttitle{\end{center}}
\preauthor{\begin{center}\vest\Large\color{frbblue}}
\postauthor{\end{center}}
\predate{\begin{center}\vest}
\postdate{\end{center}}
\usepackage{tikz}
\setlength{\parskip}{0.5\baselineskip}
\setlength{\parindent}{0cm}
\setcounter{secnumdepth}{0}
\newcommand{\qm}[1]{``#1''}

\begin{document}

\title{Engelsk B}
\author{Øvespørgsmål}
\date{}
\maketitle

\begin{tikzpicture}[remember picture,overlay]
\node[anchor=north east,yshift=-15pt,xshift=-45pt]
  at (current page.north east)
  {\includegraphics[scale=0.1]{frbvuc_logo.png}};
\end{tikzpicture}
\thispagestyle{empty}
\vspace{-4em}
\begin{center}
{\color{frbblue}\rule{\textwidth}{1pt}}
\end{center}

\subsection{Prøvemateriale}
\begin{description}
\item[Theme:] South Africa
\item[Text:] Oren Gruenbaum, \qm{`Unlock the door or we'll kick it down': why South Africa's youngest politician is in a hurry for change}. \textit{The Guardian}, 24 May 2024. Article (excerpt).

{[\ldots] angiver forkortelser i forhold til originalen.}
\item[NS:] 2,4
\item[Instruction:] You must analyse and interpret the article and make a presentation (roughly 8 minutes) of the text and relate it to the theme \qm{South Africa}.

Apart from your own points, you must include the following points in your presentation:
\begin{itemize}
\item Characterisation
\item Intention
\end{itemize}
\end{description}

\subsection{Bedømmelseskriterier, hf B}
Ved den \textbf{mundtlige prøve} lægges der vægt på, at eksaminanden
\begin{itemize}
\item behersker et sammenhængende og forholdsvis flydende engelsk med relativ høj grad af grammatisk korrekthed
\item giver en klart sammenhængende præsentation
\item analyserer, fortolker og perspektiverer prøvematerialet med anvendelse af fagets analytiske begreber og metoder
\item anvender den viden, der er opnået i arbejdet med det studerede emne.
\end{itemize}
Der lægges i bedømmelsen vægt på, at eksaminanden kan indgå i uddybende samtale om præsentationen.

Der gives én karakter ud fra en helhedsvurdering af eksaminandens præstation.

\clearpage
\subsection{`Unlock the door or we'll kick it down': why South Africa's youngest politician is in a hurry for change}
\textbf{Oren Gruenbaum}\\
\textit{The Guardian}, 24 May 2024

\begin{linenumbers}
They are called the “born-free” South Africans – too young to remember apartheid and with no instinctive loyalty to the African National Congress. And they could deliver a stinging rebuke to the party that led the liberation struggle, if this generation votes at all in elections on 29 May.

Three decades on from the birth of democracy in the country, the ANC could see its vote share fall below 50\% for the first time.

In a country that became mired in corruption under the former president Jacob Zuma, inequality has soared and a run of two months without power cuts by the state energy firm Eskom is cause for celebration.

Now a new generation of politicians, many rising from social movements, are trying to reinvigorate the ANC. Youngest of all is Fasiha Hassan, who came to prominence as a leader of \#FeesMustFall, a wave of student protests over tuition costs that swept through universities in 2015.

At 30, Hassan is the same age as democratic South Africa and in a hurry to see a transfer of power from the old guard. In 2019, the law graduate became the youngest member of the provincial legislature for Gauteng, the region around Johannesburg; now she is 65th on the ANC’s list of 200 candidates for the parliamentary elections. Young people, she says, are central to the “renewal project” of the ANC.

“We need an intergenerational mix, that expertise; but there’s no new energy or ideas to fix the country’s problems. You need a cohort of young people who are well educated, willing to roll up their sleeves and say, ‘Right, what do we do for the next 10 years?’” In a recent election video, she talks of “cleaning up the ANC” and insists: “We need a seat at the decision-making table.”

Wearing a keffiyeh draped around her shoulders, she issues a bold ultimatum: “Unlock the door or we kick it down.”

Hassan is equally impatient about the obstacles to more equal representation in South African politics. “There is a different set of rules for women and queer politicians,” she says. “A sense of ownership of young women who are public figures – that gives you very little room for mistakes. But there’s also a very large opinion on matters outside of public office: what you wear, where you go, who you’re with, who you’re not with.”

She gets criticism from “very conservative” elements in her own Asian Muslim community. “People who don’t believe a young Muslim woman should be as bold and loud. Even before I took public office, there was a lot of support but also a backlash, ‘Look at these Muslim girls, how they dress, look at them protesting.’ But we remind them that we have every right, even in Islam, to exist in this space and be leaders.”

At a recent symposium at the University of London, she spoke ruefully about older female colleagues putting her down and criticising her clothing, or being taken for a secretary or “plus-one” when travelling with male politicians – and of evading a stalker. Equality of pay and opportunity is crucial, Hassan says, but she acknowledges that gender-based violence is so pervasive in South Africa that improving safety is more important. “I’m afraid to go to my car late at night, but I have a car. It’s a huge privilege.”

To combat South Africa’s high levels of domestic violence, she has taken up innovative ideas such as empowering social workers and community activists to record statements in rape cases so survivors do not have to go immediately to the police, with all the trauma that entails, and suggests separate courts for sexual violence.

“One of the consequences of apartheid and hundreds of years of colonialism has been a very violent society. We haven’t healed,” she says. “We need a sort of Truth and Reconciliation Commission in our own communities dealing with gender violence.”

Hassan is hopeful that this election can lead to some of these ideas being taken up. “There’s a hunger in South Africa for new leadership, within the ANC and outside. They must evolve,” she says. “Evolve or die.” There is a sense of a baton being passed: “Our parents, our grandparents, fought and achieved political freedom,” she says. “But we’ve not reached a point of economic freedom. We view that as our generational mission.”

The sense of an age gap within the party also comes out in her dissenting attitude to the ANC’s old allies in Zimbabwe (“I’m not very liked by the Zanu-PF comrades there”) and Russia (“There can be no justification for invading another sovereign territory … international law is fundamental”). But she is quick to stress South Africa’s non-aligned position: “We can no longer support a unipolar world … Brics [an alliance of developing nations] can help create a world that’s more equal in terms of power dynamics; it gives us more leverage on the world stage.”

Hassan is not too young to remember more optimistic times, citing the former president Thabo Mbeki’s I am an African speech in 1996. “I remember being in primary school and feeling a deep sense of ‘we’ve a good future ahead of us,’” she says. “There was hope.” But if Mbeki turned out to be a deeply flawed leader, forced out of office by his party and tainted by his conspiratorial view of HIV/Aids, she feels his successor, Zuma, reached new lows.

“We saw the economy stagnate, but worse, we started to see the state being hollowed out,” Hassan says, referring to findings of “state capture” under Zuma. Allegations of systemic corruption keep coming. “I don’t think the party is irredeemable,” she says. “I wouldn’t stand for public office if I believed that.” But she accepts it is “an uphill battle” to fight corruption. “It’s not just the ANC – it’s all political parties. There’s a group of people who think politics is a way to make money.” It is about fixing the party, she says, “but ultimately fixing the state”.


[\ldots]
\end{linenumbers}

\end{document}
